
\documentclass[a4paper,11pt]{article}
\usepackage{fontspec}
\setmainfont{Noto Sans}
\usepackage[french]{babel}
\usepackage{microtype}
\usepackage[top=9mm,bottom=8mm,left=11mm,right=11mm]{geometry}
\usepackage{xcolor}
\usepackage{tikz}
\usetikzlibrary{arrows.meta,calc,positioning}
\usepackage[most]{tcolorbox}
\pagestyle{empty}

\definecolor{Green}{HTML}{319B58}
\definecolor{GreenDark}{HTML}{17613A}
\definecolor{GreenLight}{HTML}{EAF6EC}
\definecolor{Border}{HTML}{B9D3C3}
\definecolor{Ink}{HTML}{172A32}
\definecolor{Grey}{HTML}{64757D}
\definecolor{Red}{HTML}{D52F35}
\definecolor{RedLight}{HTML}{FCEBED}

\tcbset{
 enhanced,boxrule=.6pt,arc=2.8mm,
 left=3.7mm,right=3.7mm,top=1.75mm,bottom=1.75mm,
 boxsep=1mm,coltext=Ink
}
\newtcolorbox{greenbox}[1][]{colback=GreenLight,colframe=Border,#1}
\newtcolorbox{warnbox}[1][]{colback=RedLight,colframe=Red,
 left=2.7mm,right=2.7mm,top=1.25mm,bottom=1.25mm,#1}

\begin{document}
\begin{tikzpicture}[remember picture,overlay]
\draw[draw=Border,line width=.9pt,rounded corners=1.7mm]
 ([xshift=1mm,yshift=-1mm]current page.north west)
 -- ([xshift=-1mm,yshift=-1mm]current page.north east)
 -- ([xshift=-1mm,yshift=1mm]current page.south east)
 -- ([xshift=1mm,yshift=1mm]current page.south west)--cycle;
\fill[Green] ([xshift=1mm,yshift=-1mm]current page.north west)
 rectangle ([xshift=5.2mm,yshift=1mm]current page.south west);
\end{tikzpicture}

\vspace*{1mm}\hspace*{6mm}
{\color{GreenDark}\fontsize{25}{27}\selectfont\bfseries PRESENT DEL SUBJONTIU}

\vspace{0.5mm}\hspace*{6mm}
{\color{Grey}\large Le présent du subjonctif}

\vspace{2.5mm}\hspace*{6mm}
\begin{minipage}{0.91\textwidth}

\begin{greenbox}
\textbf{\color{GreenDark}CONSTRUCCION}\\[1mm]
\centering
{\Large \textbf{RADICAL + TERMINASONS DEL SUBJONTIU}}\\[1mm]
\small
Exemple : \textbf{MOSTRAR}\\[1mm]
\textbf{mostre}\quad \textbf{mostres}\quad \textbf{mostre}\quad
\textbf{mostrem}\quad \textbf{mostretz}\quad \textbf{mostren}
\end{greenbox}

\vspace{2.5mm}
\begin{center}
\begin{tikzpicture}[>=Latex]
\draw[Ink!65,line width=.8pt] (0,0)--(15.1,0);
\node[below=2mm] at (2.0,0) {\small PASSAT};
\node[below=2mm] at (7.5,0) {\small ARA};
\node[below=2mm] at (13.5,0) {\small FUTUR};
\draw[Green,line width=2.5pt,->] (7.7,0)--(9.8,0);
\draw[GreenDark,line width=1.1pt,->] (7.5,1.0)--(7.5,.18);
\node[above=3.8mm] at (7.5,1.0) {\bfseries\color{GreenDark}ARA};
\node[above=3mm,align=center,text width=6.2cm] at (8.7,0)
{\small\bfseries volontat, necessitat, dubte, possibilitat};
\end{tikzpicture}
\end{center}

\vspace{1.5mm}
\begin{greenbox}
\textbf{\color{GreenDark}A QUÉ SERVÍS ?}\\[1mm]
\begin{itemize}\setlength{\itemsep}{0.4mm}
\item exprimir una \textbf{volontat}, un desir o una necessitat ;
\item exprimir lo \textbf{dobte}, l'incertitud o una possibilitat ;
\item presentar una accion dependenta d'una autra proposicion ;
\item après de certanas expressions e conjuncions coma \textbf{que, cal que, vòli que, es possible que}.
\end{itemize}
\end{greenbox}

\vspace{2mm}
\begin{greenbox}
\textbf{\color{GreenDark}EXEMPLES}\\[1mm]
\textbf{Vòli que vengas amb nosautres.}\\
\textbf{Cal que trabalhem uèi.}\\
\textbf{Es possible que plòga deman.}\\
\textbf{Dubi que siá vertat.}
\end{greenbox}

\vspace{2mm}
\begin{minipage}{0.485\textwidth}
\begin{greenbox}[height=3.25cm]
\textbf{\color{GreenDark}UNA IDÈA CLAU}\\[1mm]
Lo subjontiu es sovent ligat a una proposicion que depend d'una autra :\\[1mm]
\centering
\textbf{Vòli}\\
$\downarrow$\\
\textbf{que vengas}\\[1mm]
\raggedright
La segonda accion es pas presentada coma un fach simple, mas coma una volontat.
\end{greenbox}
\end{minipage}\hfill
\begin{minipage}{0.485\textwidth}
\begin{warnbox}[height=3.25cm]
\textbf{\color{Red}ATENCION}\\[1mm]
Lo subjontiu es pas definit simplament per un mot coma « que ».\\[1mm]
Es lo \textbf{sens de la proposicion} — volontat, dobte, necessitat, possibilitat... — que compta.
\end{warnbox}
\end{minipage}

\vspace{2mm}
\begin{greenbox}
\textbf{\color{GreenDark}A RETÉNER}\\[1mm]
Lo present del subjontiu permet de presentar una accion coma \textbf{volguda, necessària,
dobtosa, possibla o dependent d'una autra proposicion}. Sa forma depend del vèrb e del grop de conjugason.
\end{greenbox}

\vspace{1.5mm}\centering
{\small\color{Grey}\textbf{Indicis frequents :} que, cal que, vòli que, dobti que, es possible que, cresi pas que...}

\vspace{1mm}
{\scriptsize\color{Grey}Referéncia : occitan languedocian, graphia classica, vèrb'Òc / Lo Congrès permanent de la lenga occitana.}

\end{minipage}
\end{document}
